\documentclass[10pt,a4paper]{article}
\usepackage[margin=21mm,headheight=15pt]{geometry}
\usepackage[T1]{fontenc}\usepackage{lmodern,amsmath,amssymb,braket,booktabs,tabularx,array,longtable,graphicx,xurl,fancyhdr}
\usepackage[colorlinks=true,linkcolor=blue,urlcolor=blue]{hyperref}
\pagestyle{fancy}\fancyhf{}\fancyhead[L]{\small Applied Sciences \textbar\ applsci-4605095}\fancyhead[R]{\small SUPPLEMENT}\fancyfoot[C]{\thepage}
\setlength{\parindent}{0pt}\setlength{\parskip}{5pt}
\newcommand{\Tr}{\operatorname{Tr}}\newcommand{\CX}{\mathrm{CX}}\newcommand{\Ccode}{\mathcal C}
\renewcommand{\thesection}{S\arabic{section}}
\renewcommand{\theequation}{S\arabic{equation}}
\renewcommand{\thetable}{S\arabic{table}}
\newcolumntype{Y}{>{\raggedright\arraybackslash}X}
\begin{document}
\begin{center}
{\LARGE\bfseries Supplementary Derivations}\par\vspace{6pt}
{\normalsize Exact Logical-Channel Equivalence and Code-Space-Leakage-Aware Tomography of a Five-Qubit Encoded-Qubit Transfer Circuit}\par\vspace{4pt}
Ami Ukai, Syunsuke Fujiwara, Takahiko Mori and Hideaki Okazaki\par
Manuscript applsci-4605095 --- revised 26 October 2026
\end{center}

This supplement expands the compact proof in Section 3 and Appendix A of the revised article. The operator identity applies to the specified ideal circuit; it does not assert exact noisy-device performance. Tensor factors are written in descending wire order. The signs and binary labels have the same meanings as in the main text.

\section{Tensor-embedded controlled-X notation}
For an $n$-qubit register $\ket{b_{n-1}\cdots b_0}$, let
\[
 O_{[j]}=I_{n-1}\otimes\cdots\otimes O_j\otimes\cdots\otimes I_0.
\]
For distinct control $i$ and target $j$,
\begin{equation}
\CX_{i,j,n}=\bigl(\ket0\!\bra0\bigr)_{[i]}+
\bigl(\ket1\!\bra1\bigr)_{[i]}X_{[j]}.
\end{equation}
The implicit identity on the target in the first term, and on every other factor in both terms, makes the embedding explicit. Equivalently, it maps $b_j$ to $b_j\oplus b_i$ and leaves every other bit unchanged. This is standard controlled-gate algebra; the extended Dirac notation is related to Yamanashi, Ozawa and Okazaki, IEICE Technical Report 122, CAS2022-61, 1--6 (2023), cited in the main article. No new controlled-gate formalism is claimed.

\section{Gate-by-gate logical-basis derivation}
Fix $x,y,m,M_1,M_2\in\{0,1\}$, write $z=m\oplus x$ and $\kappa=m\oplus M_1\oplus M_2$, and define
\begin{equation}
\ket{\lambda_L^{(y)}}=\ket{y\oplus\lambda,\lambda},\qquad
\ket{G_{m;M_1,M_2}}=\frac1{\sqrt2}\sum_{\eta=0}^1(-1)^{m\eta}
\ket{M_2\oplus\eta,M_1\oplus\eta,\eta}.
\end{equation}
Here $\lambda$ is the logical basis bit and $\eta$ is the resource branch bit. The five Hadamard summation bits below are denoted $a_1,\ldots,a_5$ and are local dummy variables. Binary exponents of $(-1)$ are evaluated modulo two; the written sums run over each indicated bit in $\{0,1\}$.

The initial five-qubit state is
\begin{equation}
\ket{\Psi^{(0)}_\lambda}=\frac1{\sqrt2}\sum_\eta(-1)^{m\eta}
\ket{M_2\oplus\eta,M_1\oplus\eta,\eta,y\oplus\lambda,\lambda}_{43210}.
\end{equation}
Read the main circuit from left to right. The first four gates $H_1,\CX_{1,2,5},H_1,\CX_{2,3,5}$ give
\begin{equation}
\ket{\Psi^{(1)}_\lambda}=\frac1{2\sqrt2}\sum_{\eta,a_1,a_2}
(-1)^{m\eta+a_1(y\oplus\lambda)+a_1a_2}
\ket{M_2\oplus\eta,M_1\oplus a_1,\eta\oplus a_1,a_2,\lambda}.
\end{equation}
The next block $H_3,X_3^m,\CX_{1,3,5},H_3,X_3^m$ yields
\begin{align}
\ket{\Psi^{(2)}_\lambda}
 &=\frac1{4\sqrt2}\sum_{\eta,a_1,a_2,a_3,a_4}(-1)^{P_2}
 \ket{M_2\oplus\eta,a_4\oplus m,\eta\oplus a_1,a_2,\lambda},\\
 P_2&=m\eta+a_1(y\oplus\lambda)+a_1a_2+a_3(M_1\oplus a_1)
       +a_4(a_3\oplus m\oplus a_2).
\end{align}
The three controlled gates $\CX_{2,4,5},\CX_{4,3,5},\CX_{0,3,5}$ change only the basis labels to
\begin{equation}
\ket{M_2\oplus a_1,a_4\oplus m\oplus M_2\oplus a_1\oplus\lambda,
\eta\oplus a_1,a_2,\lambda}.
\end{equation}
Finally, $H_4,X_4^z,\CX_{4,0,5}$ give
\begin{align}
\ket{\Psi^{(f)}_\lambda}
 &=\frac18\sum_{\eta,a_1,\ldots,a_5}(-1)^P
 \ket{a_5\oplus z,a_4\oplus m\oplus M_2\oplus a_1\oplus\lambda,
 \eta\oplus a_1,a_2,\lambda\oplus a_5\oplus z},\\
 P&=P_2+a_5(M_2\oplus a_1).
\end{align}
These expressions follow directly from
$H\ket b=2^{-1/2}\sum_a(-1)^{ba}\ket a$ and $X^c\ket b=\ket{b\oplus c}$.

Set $\zeta=\eta\oplus a_1$. The sum over $a_3$ enforces $a_4=M_1\oplus a_1$. The remaining exponent becomes
\begin{equation}
P=m\zeta+mM_1+M_1a_2+a_5M_2+a_1(y\oplus\lambda\oplus a_5).
\end{equation}
The sum over $a_1$ therefore enforces $a_5=y\oplus\lambda$. The final labels are
\begin{equation}
(q_{4'},q_{3'},q_{2'},q_{1'},q_{0'})=
(y\oplus\lambda\oplus z,\lambda\oplus\kappa,\zeta,a_2,y\oplus z).
\end{equation}
The remaining $\zeta,a_2$ sums are $H\ket m_{2'}$ and $H\ket{M_1}_{1'}$. Thus
\begin{align}
 U_T^{(x,m)}\bigl(\ket{G_{m;M_1,M_2}}\otimes\ket{\lambda_L^{(y)}}\bigr)
 &=(-1)^{mM_1+M_2(y\oplus\lambda)}
 \ket{y\oplus\lambda\oplus z,\lambda\oplus\kappa}_{4'3'}\notag\\
 &\quad\otimes H\ket m_{2'}\otimes H\ket{M_1}_{1'}\otimes\ket{m\oplus x\oplus y}_{0'}.
\end{align}
The discarded vector is independent of $\lambda$. The only $\lambda$-dependent phase is $(-1)^{M_2\lambda}$, which is exactly the $Z^{M_2}$ factor of
\begin{equation}
 C=X^z\otimes X^\kappa Z^{M_2},\qquad
 C^\dagger=X^z\otimes Z^{M_2}X^\kappa,\qquad C^\dagger C=I.
\end{equation}
The inverse is an exact operator identity, with no unspecified phase. Linearity yields the same receiver map for every encoded matrix unit $\ket{i_L}\bra{j_L}$, including $i\ne j$, and hence for arbitrary mixed states and reference-entangled inputs supported on the fixed code.

\section{Why two physical input qubits are not two logical qubits}
For a fixed parity $y$, $V_y\ket\lambda=\ket{y\oplus\lambda,\lambda}$ has a two-dimensional image. Let $P_y=V_yV_y^\dagger$. The discarded output $\ket{m\oplus x\oplus y}_{0'}$ distinguishes $y=0$ from $y=1$. Therefore, after correction, tracing out discarded wires removes the cross-parity blocks:
\begin{equation}
\mathcal R\circ\mathcal T(\rho)=P_0\rho P_0+P_1\rho P_1.
\end{equation}
The ordered input/output identification is understood. Possible parity-dependent global signs in the isometries cancel within each diagonal block and cannot restore the discarded cross-parity coherence. This proves identity only on either fixed code, not on the unrestricted four-dimensional input space. The physical measurements in run02 use $y=0$; the statement for $y=1$ is theoretical.

The supplied audit independently checks 64 basis-state instances, 128 fixed-code matrix-unit instances, and 256 unrestricted-input matrix-unit instances against the released ideal unitary. Maximum absolute discrepancies are below $10^{-15}$. These are numerical checks of the stated algebra, not extra hardware measurements.
\end{document}
